\chapter{Fibracje i długi ciąg homotopii}
\label{ch:fibracje-v4}

\section{Dlaczego samo włókno nie wystarcza?}
Wiązka nakryciowa pozwala jednoznacznie podnosić drogi. W wiązce
okręgowej podniesień jest zazwyczaj wiele, ale można je wybierać
\emph{ciągle w rodzinach}. Właśnie tej cechy potrzeba, aby porównywać
homotopię włókna, przestrzeni całkowitej i podstawy. Wszystkie
przestrzenie w tym rozdziale rozpatrujemy z ich topologiami; dla
przestrzeni odwzorowań stosujemy topologię zwarto-otwartą.

\begin{definicja}[Własność podnoszenia homotopii]
Niech $p:E\to B$ będzie ciągłe, $Z$ przestrzenią topologiczną,
$f:Z\to E$ ciągłe oraz $H:Z\times I\to B$ homotopią spełniającą
$H(z,0)=p(f(z))$. Mówimy, że $p$ ma \emph{własność podnoszenia
homotopii względem $Z$}, jeżeli istnieje ciągłe
$\widetilde H:Z\times I\to E$ o
$\widetilde H(z,0)=f(z)$ oraz $p\widetilde H=H$.
Odwzorowanie $p$ nazywamy \emph{fibracją Hurewicza}, jeśli posiada tę
własność dla każdego $Z$, oraz \emph{fibracją Serre'a}, jeśli posiada ją
dla wszystkich dysków $Z=D^k$, $k\ge0$. Włóknem nad $b\in B$ nazywamy
$F_b=p^{-1}(b)$.
\end{definicja}

\begin{przyklad}
Rzut $B\times F\to B$ jest fibracją Hurewicza: przy
$f(z)=(f_B(z),f_F(z))$ można wziąć
$\widetilde H(z,t)=(H(z,t),f_F(z))$. Nakrycie także jest fibracją
Hurewicza: lokalna jednoznaczność podnoszenia dróg umożliwia
podnoszenie całych rodzin. Włókna tej samej fibracji nad różnymi punktami podstawy nie
muszą być homeomorficzne, lecz nad łukowo spójną podstawą
fibracji Hurewicza są równoważne homotopijnie.
\end{przyklad}

\begin{lemat}[Względne podnoszenie na komórce]
\label{lem:fib-wzgledne-v4}
Niech $p:E\to B$ będzie fibracją Serre'a, $k\ge0$,
$H:D^k\times I\to B$ i
$L=(D^k\times\{0\})\cup(\partial D^k\times I)$.
Każda mapa $h:L\to E$ spełniająca $ph=H|_L$ rozszerza się do mapy
$\widetilde H:D^k\times I\to E$ spełniającej
$p\widetilde H=H$. Dla pary CW $(X,A)$ analogicznie można podnieść
homotopię na $X\times I$, jeżeli podniesienie zadano na
$(X\times\{0\})\cup(A\times I)$.
\end{lemat}
\begin{proof}
Dla $k=0$ jest to własność podnoszenia dla punktu. Gdy $k\ge1$,
$D^k\times I$ jest kulą topologiczną, a $L$ jest jej brzegiem
z usuniętym wnętrzem górnej $k$-ściany. Istnieje homeomorfizm par
$(D^k\times I,L)\cong(D^k\times I,D^k\times\{0\})$:
obie podprzestrzenie są domkniętymi $k$-dyskami w brzegu kuli;
identyfikujemy brzeg ze sferą i przesuwamy odpowiedni dysk na dolną
półsferę, po czym przedłużamy homeomorfizm radialnie do kuli.
Po przeniesieniu diagramu mamy standardową własność podnoszenia.

Dla CW-pary przeprowadzamy indukcję po szkieletach $X$.
Wchodząc na nową komórkę $e^k$ używamy jej mapy
charakterystycznej $D^k\to X$; podniesienie jest już określone
na $\partial D^k\times I$ i $D^k\times\{0\}$, więc działa pierwszy
krok. Na wspólnych brzegach otrzymane mapy zgadzają się ze
starszym podniesieniem. Topologia słaba CW, a następnie fakt,
że $I$ jest lokalnie zwarty, zapewniają ciągłość mapy sklejonej
na $X\times I$. Nie zakładamy jednoznaczności podniesienia.
\end{proof}

\begin{twierdzenie}[Długi ciąg dokładny fibracji]
\label{thm:fib-les-v4}
Niech $p:E\to B$ będzie fibracją Serre'a, $e_0\in E$,
$b_0=p(e_0)$ oraz $F=p^{-1}(b_0)$. Wówczas istnieje naturalny
ciąg dokładny grup homotopii
\[
\cdots\to\pi_n(F,e_0)\xrightarrow{i_*}\pi_n(E,e_0)
\xrightarrow{p_*}\pi_n(B,b_0)\xrightarrow{\partial}\pi_{n-1}(F,e_0)
\to\cdots.
\]
Na końcu otrzymujemy $\pi_1(B,b_0)\to\pi_0(F)\to\pi_0(E)\to\pi_0(B)$;
tu dokładność rozumie się dla zbiorów z wyróżnionym punktem.
Dla $n=2$ grupa po prawej może być nieprzemienna, choć obraz
homomorfizmu $\partial$ jest wtedy przemienny.
\end{twierdzenie}
\begin{proof}
Z długiego ciągu pary $(E,F)$ wystarczy otrzymać naturalne
izomorfizmy $\pi_n(E,F,e_0)\cong\pi_n(B,b_0)$ dla $n\ge2$
oraz odpowiednią bijekcję w stopniu $1$.
Użyjmy kostkowego opisu grup względnych: wyróżniona ściana
$I^{n-1}\times\{0\}$ jest odwzorowywana do $F$, a wszystkie
pozostałe ściany do $e_0$. Aby podnieść mapę
$(I^n,\partial I^n)\to(B,b_0)$, traktujemy jej ostatnią
współrzędną jako czas; na jednej skrajnej ścianie i na ścianach
bocznych ustalamy $e_0$. Lemat~\ref{lem:fib-wzgledne-v4}
rozszerza podniesienie. Otrzymujemy surjektywność.

Jeśli dwie mapy względne $u,v$ mają homotopijne obrazy w $B$,
podnosimy najpierw homotopię ich obrazów, przytrzymując ściany
z wartością $e_0$. Otrzymujemy $u\simeq u'$ z $pu'=pv$.
Następnie podnosimy stałą w czasie homotopię $pv$, przy
z góry ustalonych wartościach $u'$ i $v$ na obu podstawach
kostki $I^n\times I$ i na pozostałych ścianach bocznych.
Ten ustalony fragment brzegu jest dyskiem $n$-wymiarowym:
komplementem w brzegu całej kostki jest wnętrze
$\bigl(I^{n-1}\times\{0\}\bigr)\times I$. Ten sam argument
z homeomorfizmem par, co w lemacie, daje homotopię
$u'\simeq v$ przez mapy względne. Jest to iniektywność.
Zgodność z działaniem grupowym wynika z zachowania
sklejania kostek przez $p$. Wstawienie uzyskanych identyfikacji
do długiego ciągu pary daje tezę.
\end{proof}

\begin{uwaga}[Jak obliczać homomorfizm łączący?]
Dla $a:(D^n,S^{n-1})\to(B,b_0)$ podnosimy $a$ po
ustaleniu $e_0$ na odpowiedniej części brzegu. Restrykcja
podniesienia do pozostałej ściany leży w $F$ i reprezentuje
$\partial[a]\in\pi_{n-1}(F,e_0)$. Wybór orientacji kostki
może zmienić znak $\partial$, dlatego należy zachować
jedną konwencję w całym rachunku.
\end{uwaga}

\section{Transport wzdłuż drogi i monodromia}
\begin{stwierdzenie}
Niech $p:E\to B$ będzie fibracją Hurewicza, a $\gamma:I\to B$
drogą od $b_0$ do $b_1$. Istnieje transport ciągły
$T_\gamma:F_{b_0}\to F_{b_1}$ będący równoważnością
homotopijną. Jego klasa homotopii zależy tylko od klasy
$\gamma$ względem końców.
\end{stwierdzenie}
\begin{proof}
Stosujemy własność podnoszenia do $Z=F_{b_0}$,
$f:F_{b_0}\hookrightarrow E$ i $H(x,t)=\gamma(t)$, a następnie
kładziemy $T_\gamma(x)=\widetilde H(x,1)$.
Podobnie wybieramy transport wzdłuż drogi odwróconej.
Aby wykazać, że są to odwrotności z dokładnością do homotopii,
zauważmy, że $\gamma*\overline\gamma$ jest homotopijna względem
końców drodze stałej. Złączone podniesienia dają transport wzdłuż
złożenia; własność podnoszenia homotopii rodzin dróg zapewnia
homotopię z transportem wzdłuż drogi stałej. Ten ostatni można
wybrać jako identyczność. Analogicznie dowodzimy
$T_\gamma T_{\overline\gamma}\simeq\id$.
\end{proof}

\begin{przyklad}[Butelka Kleina]
Dla $K=(S^1\times I)/((z,1)\sim(\overline z,0))$
rzut na $I/(0\sim1)\cong S^1$ jest wiązką o włóknie $S^1$.
Monodromia działa na $\pi_1(S^1)=\Z$ przez $m\mapsto-m$.
Z tego wynika prezentacja
$\pi_1(K)\cong\langle a,b\mid aba^{-1}=b^{-1}\rangle$.
Grupa nie jest iloczynem prostym dwóch kopii $\Z$.
\end{przyklad}

\section{Fibracja przestrzeni dróg}
Dla punktowanej przestrzeni $(B,b_0)$ określmy
$PB=\{\gamma:I\to B:\gamma(0)=b_0\}$ oraz
$\Omega B=\{\gamma\in PB:\gamma(1)=b_0\}$.
\begin{twierdzenie}
Ewaluacja $\mathrm{ev}_1:PB\to B$, $\gamma\mapsto\gamma(1)$,
jest fibracją Hurewicza z włóknem $\Omega B$, a $PB$ jest
ściągalna. Dla $n\ge2$ zachodzi
$\pi_n(B,b_0)\cong\pi_{n-1}(\Omega B,c_{b_0})$;
ponadto $\pi_1(B,b_0)$ jest w naturalnej bijekcji z
$\pi_0(\Omega B)$.
\end{twierdzenie}
\begin{proof}
Niech $f:Z\to PB$ oraz $H:Z\times I\to B$ spełniają
$H(z,0)=f(z)(1)$. Wybieramy podniesienie $\widetilde H(z,t)$
łącząc najpierw $f(z)$, a następnie przebiegając odcinek
$H(z,s)$ dla $0\le s\le t$. Żeby zachować warunek
$\widetilde H(z,0)=f(z)$, stosujemy zmienny podział parametru:
\[
\widetilde H(z,t)(u)=
\begin{cases}
f(z)\!\left(\dfrac{u}{1-t/2}\right),&0\le u\le1-t/2,\\[4pt]
H\!\left(z,2u-2+t\right),&1-t/2\le u\le1.
\end{cases}
\]
Przy $t=0$ drugi przedział jest jednopunktowy; obie gałęzie
mają zgodną wartość na styku, a koniec drogi to $H(z,t)$.
Lemat o sklejaniu i własność wykładnicza topologii
zwarto-otwartej dla zwartego, lokalnie zwartego $I$ dają
ciągłość całej rodziny. Homotopia
$R_s(\gamma)(u)=\gamma((1-s)u)$ ściąga $PB$ do drogi stałej.
Teraz stosujemy Twierdzenie~\ref{thm:fib-les-v4}.
\end{proof}

\section{Fibracja Hopfa: poprawny wzór i obliczenie}
\begin{twierdzenie}
\label{thm:hopf-v4}
Odwzorowanie
$h:S^3\to S^2\subset\C\times\R$ dane wzorem
\[
h(z,w)=(2z\overline w,\ |z|^2-|w|^2),\qquad
|z|^2+|w|^2=1,
\]
jest lokalnie trywialną wiązką o włóknie $S^1$.
W konsekwencji $\pi_3(S^2)\cong\Z$, a $[h]$ jest generatorem.
\end{twierdzenie}
\begin{proof}
Równość $4|z\overline w|^2+(|z|^2-|w|^2)^2=1$
pokazuje, że obraz należy do $S^2$.
Dla $\lambda\in S^1$ mamy
$h(\lambda z,\lambda w)=h(z,w)$, ponieważ
$\lambda\overline\lambda=1$.
Dlatego $h$ przechodzi na $\C P^1$.
Dla $z\ne0$ współrzędna $u=w/z$ jednoznacznie opisuje
prostą zespoloną, której jednostkowy przedstawiciel to
$s(u)=(1,u)/\sqrt{1+|u|^2}$. Mapa
$(u,\lambda)\mapsto\lambda s(u)$ jest homeomorfizmem
lokalnej trywializacji. Analogicznie postępujemy dla $w\ne0$.
W tych współrzędnych $h$ daje standardową identyfikację
$\C P^1\cong S^2$, więc jest wiązką $S^1\to S^3\to S^2$.
Z ciągu dokładnego oraz $\pi_k(S^1)=0$ dla $k\ge2$ dostajemy
$\pi_3(S^3)\xrightarrow{\cong}\pi_3(S^2)$.
Ponieważ $S^3$ jest $2$-spójna, twierdzenie Hurewicza daje
$\pi_3(S^3)\cong H_3(S^3;\Z)=\Z$; obraz klasy identyczności
jest klasą $h$.
\end{proof}
\begin{uwaga}
Zapis $2zw$ bez sprzężenia jest \emph{błędny}:
po pomnożeniu obu współrzędnych przez $\lambda\in S^1$
pierwsza współrzędna zmieniłaby się przez czynnik $\lambda^2$.
Włókna nie byłyby zatem orbitami działania $S^1$.
\end{uwaga}

\section{Zastąpienie odwzorowania fibracją}
\begin{twierdzenie}
Niech $f:X\to B$ będzie ciągłe. Zdefiniujmy
$E_f=\{(x,\gamma)\in X\times B^I:\gamma(0)=f(x)\}$,
$q(x,\gamma)=\gamma(1)$ i $j(x)=(x,c_{f(x)})$.
Wtedy $q:E_f\to B$ jest fibracją Hurewicza,
$j:X\to E_f$ jest równoważnością homotopijną oraz $qj=f$.
\end{twierdzenie}
\begin{proof}
Podnosząc homotopię w $B$, dopisujemy do drugiej współrzędnej
$(x,\gamma)$ odpowiedni odcinek drogi; zmienną parametryzację
wybieramy jak w dowodzie dla $PB$. Pierwsza współrzędna jest
niezmienna, więc $\gamma(0)=f(x)$ pozostaje prawdziwe.
Rzut $r:E_f\to X$ jest lewą odwrotnością $j$,
a $K_s(x,\gamma)=(x,\gamma((1-s)\,\cdot))$ jest
homotopią od $\id_{E_f}$ do $jr$.
\end{proof}

Włóknem homotopijnym $f$ nad $b_0$ nazywamy
$\operatorname{hofib}_{b_0}(f)=q^{-1}(b_0)$.
Dla każdej mapy można więc zastosować ciąg dokładny
\emph{włókno homotopijne -- dziedzina -- przeciwdziedzina}.

\section{Granice metody i zapowiedź ciągu Serre'a}
Długi ciąg dokładny mierzy grupy homotopii, ale nie wyznacza
sam z siebie pierścienia kohomologii $E$. Dla fibracji Serre'a
$F\to E\to B$ o podstawie łukowo spójnej i przestrzeniach
odpowiedniego typu CW ciąg spektralny Serre'a w homologii
ma drugą stronę
$E^2_{p,q}=H_p(B;\mathcal H_q(F;\Z))$,
gdzie $\mathcal H_q(F;\Z)$ jest systemem lokalnym wyznaczonym
przez monodromię. Jego zbieżność oznacza filtrację $H_{p+q}(E)$,
a nie bezpośrednie utożsamienie z sumą wyrazów strony drugiej.
Gdy monodromia jest trywialna, lokalne współczynniki można
zastąpić stałymi. Nie stosujemy tu tego ciągu bez osobnego
dowodu jego konstrukcji i warunków zbieżności.
\clearpage

